\documentclass[a4paper,10pt]{article}
\newcommand\subdir{/home/koen/LaTeX-setup/}
\input{\subdir/LaTeX-files/imports.tex}
\input{\subdir/LaTeX-files/master-internship.tex}
\usepackage{\subdir/LaTeX-files/aastex}
\usepackage{natbib}
\bibliographystyle{\subdir/LaTeX-files/astroadstitle.bst}

%Set the title, Author and date
\title{Notes Week 33}
\author{Koen Essers}
\tutorial{Master Internship}
\date{\today}

%Set the header style
\header{\tutorialstring}

%Begin the document
\begin{document}
\setuplayout
\section{Envelope mass}

\subsection{Introduction}
Last week, I showed a couple of plots in which RLOF occured during the interpulse period of the TPAGB
star.
When the mass ratios were not extreme, and not a lot of mass was lost to a circumbinary disk,
RLOF appeared to be very stable\sidenote{In the sense that the radius of the star never became larger than the
	Roche lobe radius.
}. For example, see this figure:

\begin{figure}[ht]
	\centering
	\incplot{w32-no-RLOF}
\end{figure}

Maybe this surprising stability is due to the low envelope mass at the moment of RLOF.
Below, I plot an approximate mass for the start of RLOF.
It is approximate because it is difficult to
select an exact starting moment.
This is also why there are differences in the last two columns.

\begin{figure}[ht]
	\marginfignote{6}{The third column shows models that interact during the last interpulse phase,
		where the radius of the star never exceeds the Roche lobe radius.
		As we can see, the envelope mass
		at this point is already quite small.
		But really, to know if this is really the reason, we should compare
		if with stars with different masses, but equally small evelope masses before the final thermal pulse.
	}
	\centering
	\incplot{w32-pcolor-5-m1.4}
\end{figure}

\subsection{Setup}

Let me plot the masses of the envelope for the different thermal pulses as a function of the mass.

\begin{figure}[ht]
	\marginfignote{0}{The dots represent the mass of the TPAGB star at the start of a certain thermal pulse.
		The exact thermal pulse numbers are irrelevant here.
		It is only important to look at the lowest two dot.
		These dot represent the mass of the envelope during the \emph{last} two thermal pulse.
		Clearly, if we interact during
		the last \emph{interpulse} phase, RLOF will occur when the star has an envelope mass between these two values.
	}
	\marginfignote{17}{The figure above shows that for the models that interact during the interpulse, the envelope mass is approximately \(0.5\; M_\odot \).
		This figure shows that this is totally in agreement with the values that we expect are possible.
		The envelope mass of the \(1.4\; M_\odot \) star during the last thermal pulse is not very low compared to the other stars.
	}
	\centering
	\incplot{w33-mass-at-TP}
\end{figure}

There are some quite heavy stars that have \emph{lower} envelope masses towards the end
of the last interpulse period.
For example, it might be interesting
to look at the \(2.4\; M_\odot \) model.
This one has an envelope mass during the
last thermal pulse close to \(0.1\; M_\odot \), about \(2 \times \) lower than the \(1.4\; M_\odot \) model.

I will use the simple binary code to estimate a nice Roche lobe radius that hits causes RLOF to happen towards the end of the last interpulse, when there
is very little mass left.
It is then interesting to see if we again obtain these stable RLOF mass transfer binary systems.

\begin{figure}[ht]
	\marginfignote{0}{I ended up selecting a single initial Roche lobe radius of \(R_\text{RL} = 1100\; R_\odot \).
		I simulate the same 5 initial mass ratios as I do for my
		large grid, and the same 3 \(f_{\delta } \) values, however this time only for \(\varepsilon = 0\), which means \(f_{\delta } = \delta \).
	}
	\marginfignote{10}{This seems like it will give interesting results.
		The models \emph{should} undergo RLOF with very little envelope mass left, or in the case of the \(q=1\) model,
		experience so sort of flipping where now, wind mass-loss causes an increase in the Roche lobe radius.
		This is probably some non-trivial effect of the changing mass ratio in the prescription from \citet{saladino2018} and a nice consistency check.
		I expect roughly the same behaviour in the detailed \emph{MESA} models.
	}
	\centering
	\incplot{w33-test-higher-mass-low-env}
\end{figure}

I expect there to be \emph{big} differences in the final periods \emph{and} the final [s/Fe] values for these models.
This is simply because this heavier star reaches much larger radii and much larger envelope s-process abundances.

\subsection{Results}

\newsection{Mixing mass and the Effect on the final abundances}

Last week, I showed this plot:
\begin{figure}[ht]
	\marginfignote{1}{Every plot here makes some intuitive sense if we consider that the envelope
		s-process abundances increase after thermal pulses up to the point were we start actually \emph{accreting} material
		onto a main sequence star, and we obtain MS abundances.}
	\marginfignote{65}{This is where things start to become a little difficult to predict.
		This panel shows the abundances in a MS accretor if the MS itself has no initial s-process abundances.
		To understand these abundances, we need to consider both the mixing mass, \emph{and} the accreted mass.
	}
	\marginfignote{83}{For the actual main sequence abundances, it is important to also consider the initial abundances of the s-process elements in the main sequence star.}
	\centering
	\incplot{w32-radius-4}
\end{figure}

Let me plot the \(\mu\)-mixing profiles for the blue models (\(f_{\delta } = 0.1\)).
That way, we can see
how deep the thermohaline mixing should extend, and what the differences are between the two.

\begin{figure}[ht]
	\marginfignote{0}{The orange line has less accretion, but with higher s-process abundances, and thus a higher mean molecular weight.
		The lower accretion mass is more important in this case.
		The \(\varepsilon =0.0\) model has a lower mixing mass.
	}
	\centering
	\incplot{w33-mixing-mass-1}
\end{figure}

To confirm the mixing mass actually fixes MS abundances, let me show, explicitly,
a calculation of the final accretor abundances.
First, let's consider the
case \emph{without} any initial abundances in the accretor.
Then we simply have

\begin{align*}
	X_\text{MS}                       & = \Delta M_\text{s-process} / (M_\text{mix} + \Delta M_\text{accr})        \\
	X_\text{MS, \(\varepsilon =0.1\)} & \approx 9 \times  10^{-9} / (0.09 + 0.464) \approx 1.62 \times  10^{-8}    \\
	X_\text{MS, \(\varepsilon =0.0\)} & \approx 5.5 \times  10^{-9} / (0.015 + 0.359) \approx 1.47 \times  10^{-8}
\end{align*}
These values, although approximate, align very well with what we see in the subplot showing \(\langle X \rangle (s\text{, MS} )_{t=t}, X(s,i) = 0 \)
Now, if we consider the initial abundances of the accretor, I calculate the abundances as follows:
\begin{align*}
	X_\text{MS}                       & = (\Delta M_\text{s-process} + X(s) \cdot M_\text{mix}) / (M_\text{mix} + \Delta M_\text{accr})              \\
	X_\text{MS, \(\varepsilon =0.1\)} & \approx (9 \times  10^{-9} + 4.4 \times 10^{-8} \cdot 0.464 ) / (0.09 + 0.464) \approx 5.31 \times  10^{-8}  \\
	X_\text{MS, \(\varepsilon =0.0\)} & \approx (5.5 \times  10^{-9}+ 4.4 \times 10^{-8} \cdot 0.359) / (0.015 + 0.359) \approx 5.69 \times  10^{-8}
\end{align*}
Both of these values align very well with the plot showing \(\langle X \rangle (s\text{, MS} )_{t=t},\)\\\(X(s,i) = X(s)_\text{Asplund}\).

\bibliography{master.bib}
\end{document}
